Analog radio-over-fiber (A-RoF) provides a means of transporting radio signals between centralized processing and remote radio units while retaining an analog optical interface. At millimeter-wave frequencies, its optical and wireless segments must be considered together because the recovered waveform depends on the cascaded component responses. Modulation and amplification can introduce nonlinear distortion, while filtering, fiber propagation, and receiver noise also affect signal quality. Consequently, equalizer performance should be assessed using a physical model that distinguishes linear channel effects from nonlinear component behavior and identifies the assumptions underlying its noise and power budget~\cite{song2024performance}.

A feed-forward equalizer (FFE) trained using the least-mean-square (LMS) algorithm offers a useful adaptive linear baseline. Its finite impulse response can compensate linear distortion within the available observation window, and its coefficient updates can be implemented using regular multiply--accumulate operations. However, a linear mapping cannot in general reproduce the inverse of a nonlinear component characteristic. This limitation motivates the investigation of neural equalizers when amplitude compression or modulation nonlinearities produce signal-dependent distortion. Any resulting benefit must be established against a suitably tuned baseline under the same physical conditions and with independent evaluation data.

Previous work has demonstrated the relevance of neural nonlinear compensation in millimeter-wave RoF systems. Liu \emph{et al.}\ investigated a complex-valued multilevel ANN equalizer for 16-QAM and 64-QAM transmission in a 60-GHz RoF link, addressing cross-modulation between the in-phase and quadrature components~\cite{liu2017multilevel}. In a separate 60-GHz experiment using BPSK, Liu \emph{et al.}\ studied an adaptive activation function and compared its convergence and equalization performance with conventional nonlinear compensation~\cite{liu2017adaptive}. These studies motivate examining both the nonlinear receiver mapping and its activation function. Their experimental configurations differ, so their numerical results cannot be transferred directly to the link considered here.

Building on this motivation, the present study considers neural equalization together with the arithmetic and control required for online training. In a hardware implementation, the activation function determines both the forward evaluation and the derivative used during weight adaptation. Network dimensions, numerical precision, rounding, saturation, and update ordering therefore affect the implemented learning process as well as the inference datapath. The research objective is to identify useful BER or EVM improvements over LMS equalization and then assess whether the selected neural configurations can preserve those improvements within a consistent Q8.24 training and inference formulation. The initial implementation target is a sequential FPGA engine operating at a controlled sample rate; processing the full communication waveform at its symbol rate is a separate architectural objective.
